%%%%%%%%%%%%%%%%%%%%%%%%%%%%%%%%%%%%%%%%%%%%%%%%%%%%%%%%%%%%%
%% THÉORIE DES CATÉGORIES — Chapitres 7–8
%% Catégories Abéliennes et Suites Exactes /
%% Catégories Dérivées — Introduction
%%%%%%%%%%%%%%%%%%%%%%%%%%%%%%%%%%%%%%%%%%%%%%%%%%%%%%%%%%%%%

%%%%%%%%%%%%%%%%%%%%%%%%%%%%%%%%%%%%%%%%%%%%%%%%%%%%%%%%%%%%
%%           CHAPITRE 7 : CATÉGORIES ABÉLIENNES           %%
%%                  ET SUITES EXACTES                      %%
%%%%%%%%%%%%%%%%%%%%%%%%%%%%%%%%%%%%%%%%%%%%%%%%%%%%%%%%%%%%

\chapter{Catégories abéliennes et suites exactes}\label{ch:abeliennes}
\minitoc

\vspace{1em}

Les catégories abéliennes\index{catégorie!abélienne|(} constituent le cadre
algébrique dans lequel se développe l'algèbre homologique. Introduites par
Grothendieck\index{Grothendieck, Alexander} dans son célèbre article
\emph{Tôhoku}\index{Tôhoku (article)} (1957), elles axiomatisent les
propriétés essentielles des catégories de modules et de faisceaux qui
permettent de définir les foncteurs dérivés et de manipuler les suites
exactes. Dans ce chapitre, nous construisons la théorie pas à pas : catégories
pré-additives, additives, puis abéliennes, avant d'énoncer et de démontrer
les lemmes fondamentaux (serpent, cinq) qui sont les outils de base
du calcul homologique.

\index{suite exacte|(}

%% ============================================================
\section{Catégories pré-additives}\label{sec:preadditive}
%% ============================================================

\begin{definition}[Catégorie pré-additive]\label{def:preadditive}
\index{catégorie!pré-additive}
\index{enrichissement!sur $\Ab$}
Une catégorie $\mathcal{A}$ est dite \textbf{pré-additive}
(ou \textbf{$\Ab$-catégorie}\index{Ab-catégorie@$\Ab$-catégorie})
si, pour tous objets $A, B \in \Ob(\mathcal{A})$,
l'ensemble $\Hom_{\mathcal{A}}(A, B)$ est muni d'une structure de
groupe abélien\index{groupe abélien} telle que la composition
\[
  \circ \colon \Hom_{\mathcal{A}}(B, C) \times \Hom_{\mathcal{A}}(A, B)
  \longrightarrow \Hom_{\mathcal{A}}(A, C)
\]
est bilinéaire\index{bilinéarité de la composition} :
pour tous $f, f' \colon A \to B$ et $g, g' \colon B \to C$,
\[
  g \circ (f + f') = g \circ f + g \circ f', \qquad
  (g + g') \circ f = g \circ f + g' \circ f.
\]
\end{definition}

\begin{exemple}\label{ex:preadditive}
\index{catégorie!des modules}
\begin{enumerate}[label=(\roman*)]
\item Pour tout anneau $R$, la catégorie $R\text{-}\Mod$\index{R-Mod@$R\text{-}\Mod$}
  est pré-additive : si $f, g \colon M \to N$ sont des $R$-morphismes,
  on pose $(f + g)(m) = f(m) + g(m)$.

\item La catégorie $\Ab$\index{Ab@$\Ab$} des groupes abéliens
  est le cas $R = \Z$.

\item La catégorie $\Set$\index{Set@$\Set$} n'est \emph{pas}
  pré-additive, car $\Hom_{\Set}(A, B)$ n'a pas de structure naturelle
  de groupe abélien (sauf cas dégénérés).
\end{enumerate}
\end{exemple}

\begin{remarque}\label{rem:zero-morphism}
\index{morphisme nul}
Dans une catégorie pré-additive, chaque $\Hom(A, B)$ possède un
élément neutre $0_{A,B}$ pour l'addition. On a
$0_{B,C} \circ f = 0_{A,C}$ et $g \circ 0_{A,B} = 0_{A,C}$
pour tous $f \colon A \to B$ et $g \colon B \to C$.
\end{remarque}

\begin{definition}[Foncteur additif]\label{def:additive-functor}
\index{foncteur!additif}
Un foncteur $F \colon \mathcal{A} \to \mathcal{B}$ entre catégories
pré-additives est \textbf{additif} si, pour tous objets $A, B$ et tous
morphismes $f, g \colon A \to B$,
\[
  F(f + g) = F(f) + F(g).
\]
Autrement dit, $F_{A,B} \colon \Hom_{\mathcal{A}}(A, B) \to
\Hom_{\mathcal{B}}(FA, FB)$ est un morphisme de groupes abéliens.
\end{definition}


%% ============================================================
\section{Catégories additives et biproduits}\label{sec:additive}
%% ============================================================

\begin{definition}[Objet nul]\label{def:zero-object}
\index{objet nul}
\index{objet initial}
\index{objet terminal}
Un objet $0 \in \Ob(\mathcal{A})$ est un \textbf{objet nul}
s'il est à la fois initial et terminal, c'est-à-dire si
$\Hom(0, A)$ et $\Hom(A, 0)$ sont des singletons pour tout objet $A$.
\end{definition}

\begin{definition}[Biproduit]\label{def:biproduct}
\index{biproduit}
\index{somme directe!catégorique}
Soit $\mathcal{A}$ une catégorie pré-additive possédant un objet nul.
Un \textbf{biproduit}\index{biproduit|textbf} de deux objets $A, B$ est un
diagramme
\[
\begin{tikzcd}
A \arrow[r, "\iota_A"] & A \oplus B
  \arrow[l, bend right, "\pi_A"']
  \arrow[r, bend left, "\pi_B"]
& B \arrow[l, "\iota_B"']
\end{tikzcd}
\]
vérifiant :
\begin{enumerate}[label=(\alph*)]
\item $\pi_A \circ \iota_A = \id_A$, \quad $\pi_B \circ \iota_B = \id_B$;
\item $\pi_B \circ \iota_A = 0$, \quad $\pi_A \circ \iota_B = 0$;
\item $\iota_A \circ \pi_A + \iota_B \circ \pi_B = \id_{A \oplus B}$.
\end{enumerate}
\end{definition}

\begin{proposition}\label{prop:biproduct-product-coproduct}
\index{biproduit!produit et coproduit}
Dans une catégorie pré-additive, un biproduit $A \oplus B$ est
simultanément un produit $(A \oplus B, \pi_A, \pi_B)$ et un
coproduit $(A \oplus B, \iota_A, \iota_B)$.
\end{proposition}

\begin{proof}
Montrons que $A \oplus B$ est un produit. Soit $X$ un objet muni de
morphismes $f \colon X \to A$ et $g \colon X \to B$. Posons
$h = \iota_A \circ f + \iota_B \circ g \colon X \to A \oplus B$.
Alors
\[
  \pi_A \circ h = \pi_A \circ \iota_A \circ f + \pi_A \circ \iota_B \circ g
  = f + 0 = f,
\]
et de même $\pi_B \circ h = g$. Pour l'unicité, si $h'$ vérifie
$\pi_A \circ h' = f$ et $\pi_B \circ h' = g$, alors
\[
  h' = (\iota_A \circ \pi_A + \iota_B \circ \pi_B) \circ h'
  = \iota_A \circ f + \iota_B \circ g = h.
\]
L'argument dual montre que c'est un coproduit.
\end{proof}

\begin{definition}[Catégorie additive]\label{def:additive-cat}
\index{catégorie!additive}
Une catégorie $\mathcal{A}$ est \textbf{additive} si :
\begin{enumerate}[label=(\roman*)]
\item $\mathcal{A}$ est pré-additive ;
\item $\mathcal{A}$ possède un objet nul ;
\item tout couple d'objets $(A, B)$ admet un biproduit $A \oplus B$.
\end{enumerate}
\end{definition}

\begin{exemple}\label{ex:additive-cats}
\index{catégorie!des espaces vectoriels}
\begin{enumerate}[label=(\roman*)]
\item $R\text{-}\Mod$ est additive pour tout anneau $R$, avec
  $M \oplus N$ la somme directe usuelle.
\item $\Ab$ est additive.
\item La catégorie des espaces vectoriels de dimension finie
  $\mathsf{Vect}_k^{\mathrm{fd}}$\index{Vect@$\mathsf{Vect}_k$}
  est additive.
\end{enumerate}
\end{exemple}

\begin{proposition}[Matrices de morphismes]\label{prop:matrix-morphisms}
\index{matrice de morphismes}
Dans une catégorie additive, un morphisme
$f \colon A_1 \oplus A_2 \to B_1 \oplus B_2$ est entièrement déterminé
par la matrice
\[
  f \longleftrightarrow
  \begin{pmatrix} f_{11} & f_{12} \\ f_{21} & f_{22} \end{pmatrix},
  \qquad f_{ij} = \pi_{B_i} \circ f \circ \iota_{A_j}
  \colon A_j \to B_i,
\]
et la composition correspond au produit matriciel.
\end{proposition}

\begin{proof}
On a $f = \sum_{i,j} \iota_{B_i} \circ f_{ij} \circ \pi_{A_j}$
par la condition (c) du biproduit appliquée au domaine et au codomaine.
La vérification du produit matriciel pour la composition est un calcul
direct utilisant les relations d'orthogonalité.
\end{proof}


%% ============================================================
\section{Catégories abéliennes}\label{sec:abelian}
%% ============================================================

\begin{definition}[Noyau et conoyau]\label{def:kernel-cokernel}
\index{noyau!catégorique}
\index{conoyau}
Soit $f \colon A \to B$ un morphisme dans une catégorie pré-additive
avec objet nul.
\begin{enumerate}[label=(\roman*)]
\item Un \textbf{noyau}\index{noyau|textbf} de $f$ est un
  égalisateur\index{égalisateur} de $f$ et $0_{A,B}$, c'est-à-dire un
  morphisme $\ker f \colon K \to A$ tel que $f \circ \ker f = 0$ et,
  pour tout $g \colon X \to A$ avec $f \circ g = 0$, il existe un unique
  $\bar{g} \colon X \to K$ avec $\ker f \circ \bar{g} = g$ :
\[
\begin{tikzcd}
X \arrow[dr, "g"'] \arrow[r, dashed, "\exists!\,\bar{g}"]
& K \arrow[d, "\ker f"] \\
& A \arrow[r, "f"] & B
\end{tikzcd}
\]

\item Un \textbf{conoyau}\index{conoyau|textbf} de $f$ est un
  coégalisateur\index{coégalisateur} de $f$ et $0_{A,B}$, c'est-à-dire un
  morphisme $\coker f \colon B \to C$ tel que $\coker f \circ f = 0$ et
  vérifiant la propriété universelle duale.
\end{enumerate}
\end{definition}

\begin{definition}[Image et coimage]\label{def:image-coimage}
\index{image!catégorique}
\index{coimage}
Soit $f \colon A \to B$ un morphisme dans une catégorie possédant
noyaux et conoyaux.
\begin{enumerate}[label=(\roman*)]
\item L'\textbf{image} de $f$ est $\im f = \ker(\coker f)$.
\item La \textbf{coimage} de $f$ est $\operatorname{coim} f = \coker(\ker f)$.
\end{enumerate}
\end{definition}

\begin{definition}[Catégorie abélienne]\label{def:abelian-cat}
\index{catégorie!abélienne|textbf}
\index{axiomes de Grothendieck}
Une catégorie $\mathcal{A}$ est \textbf{abélienne} si :
\begin{enumerate}[label=(AB\arabic*)]
\item\label{AB0} $\mathcal{A}$ est additive ;
\item\label{AB1} tout morphisme admet un noyau et un conoyau ;
\item\label{AB2} pour tout morphisme $f \colon A \to B$, le morphisme
  canonique
  \[
    \operatorname{coim} f \longrightarrow \im f
  \]
  est un isomorphisme.
\end{enumerate}
\end{definition}

\begin{remarque}\label{rem:abelian-reformulation}
\index{monomorphisme!dans une catégorie abélienne}
\index{épimorphisme!dans une catégorie abélienne}
La condition \ref{AB2} équivaut à : tout monomorphisme est un noyau et tout
épimorphisme est un conoyau. Plus précisément, dans une catégorie
abélienne, un morphisme $f$ se factorise comme
\[
\begin{tikzcd}
A \arrow[r, two heads, "\operatorname{coim} f"]
& \im f \arrow[r, hook, "\im f"]
& B
\end{tikzcd}
\]
où la première flèche est un épimorphisme et la seconde un monomorphisme.
C'est l'analogue catégorique du premier théorème d'isomorphisme\index{théorème d'isomorphisme!premier}.
\end{remarque}

\begin{theoreme}[Théorème de plongement de Freyd--Mitchell]%
\label{thm:freyd-mitchell}
\index{théorème!de Freyd--Mitchell}
\index{Freyd, Peter}
\index{Mitchell, Barry}
Toute petite catégorie abélienne $\mathcal{A}$ admet un foncteur
exact et pleinement fidèle
\[
  F \colon \mathcal{A} \longrightarrow R\text{-}\Mod
\]
pour un certain anneau $R$. En particulier, toute relation
diagrammatique\index{chasse au diagramme} valable dans les catégories
de modules est valable dans toute catégorie abélienne.
\end{theoreme}

\begin{proof}
La démonstration complète, due à Freyd (1964) et Mitchell (1965),
est longue et technique. L'idée est la suivante :
\begin{enumerate}[label=\arabic*.]
\item On fixe un générateur $U$ de $\mathcal{A}$
  (qui existe car $\mathcal{A}$ est petite abélienne).
\item On pose $R = \End_{\mathcal{A}}(U)^{\op}$ et on considère le
  foncteur $F = \Hom_{\mathcal{A}}(U, -) \colon \mathcal{A} \to R\text{-}\Mod$.
\item On montre que $F$ est exact et pleinement fidèle en utilisant
  le fait que $U$ est un générateur.
\end{enumerate}
Nous renvoyons à \cite{freyd1964} ou \cite{mitchell1965} pour les
détails.
\end{proof}

\begin{remarque}\label{rem:freyd-mitchell-usage}
Le théorème de Freyd--Mitchell\index{théorème!de Freyd--Mitchell}
justifie la technique de la \emph{chasse au diagramme}\index{chasse au diagramme|textbf} :
pour démontrer un résultat dans une catégorie abélienne abstraite,
il suffit de le vérifier pour les modules, où l'on peut travailler
avec des éléments.
\end{remarque}


%% ============================================================
\section{Exemples de catégories abéliennes}\label{sec:abelian-examples}
%% ============================================================

\begin{exemple}[Modules sur un anneau]\label{ex:R-mod-abelian}
\index{R-Mod@$R\text{-}\Mod$!catégorie abélienne}
Pour tout anneau $R$ (commutatif ou non), la catégorie
$R\text{-}\Mod$ est abélienne. Les noyaux et conoyaux sont les
noyaux et conoyaux de morphismes de modules au sens usuel,
et la condition image $=$ coimage est le premier théorème
d'isomorphisme : $M/\ker f \cong \im f$.
\end{exemple}

\begin{exemple}[Groupes abéliens]\label{ex:ab-abelian}
\index{Ab@$\Ab$!catégorie abélienne}
$\Ab = \Z\text{-}\Mod$ est le cas particulier $R = \Z$.
\end{exemple}

\begin{exemple}[Faisceaux de modules]\label{ex:sheaves-abelian}
\index{faisceau!de modules}
\index{catégorie!des faisceaux}
Soit $(X, \mathcal{O}_X)$ un espace annelé\index{espace annelé}.
La catégorie $\mathcal{O}_X\text{-}\mathsf{Mod}$ des faisceaux de
$\mathcal{O}_X$-modules est abélienne. Les noyaux sont calculés
section par section\index{noyau!de faisceaux} (le noyau préfaisceautique
est déjà un faisceau), mais les conoyaux\index{conoyau!de faisceaux}
nécessitent une faisceautisation.
\end{exemple}

\begin{exemple}[Foncteurs vers $\Ab$]\label{ex:functor-category-abelian}
\index{catégorie!de foncteurs!abélienne}
Si $\mathcal{C}$ est une petite catégorie, alors la catégorie des foncteurs
$\Fun(\mathcal{C}, \Ab)$\index{Fun@$\Fun(\mathcal{C}, \Ab)$} est abélienne,
les noyaux et conoyaux étant calculés argument par argument
(« point par point »\index{point par point}).
\end{exemple}

\begin{exemple}[Contre-exemple : groupes non abéliens]\label{ex:grp-not-abelian}
\index{Grp@$\Grp$!pas abélienne}
La catégorie $\Grp$ des groupes (non nécessairement abéliens)
n'est \emph{pas} abélienne. Elle est pointée (le groupe trivial
est un objet nul) et admet noyaux et conoyaux, mais la condition
\ref{AB2} échoue : le morphisme canonique
$\operatorname{coim} f \to \im f$ n'est pas toujours un isomorphisme.
En effet, $\im f$ au sens catégorique est le plus petit sous-groupe
normal contenant l'image ensembliste, qui peut différer du quotient
$G/\ker f$.
\end{exemple}


%% ============================================================
\section{Suites exactes}\label{sec:exact-sequences}
%% ============================================================

\begin{definition}[Suite exacte en un point]\label{def:exact-at}
\index{suite exacte!en un point}
Soit $\mathcal{A}$ une catégorie abélienne. Une suite de morphismes
\[
  A \xrightarrow{f} B \xrightarrow{g} C
\]
est dite \textbf{exacte en $B$} si $\im f = \ker g$
(comme sous-objets\index{sous-objet} de $B$).
De manière équivalente, $g \circ f = 0$ et, pour tout $h \colon X \to B$
avec $g \circ h = 0$, le morphisme $h$ se factorise (uniquement) par
$\im f$.
\end{definition}

\begin{definition}[Suite exacte courte]\label{def:ses}
\index{suite exacte!courte}
\index{SES|see{suite exacte courte}}
Une \textbf{suite exacte courte} est une suite
\[
\begin{tikzcd}
0 \arrow[r] & A \arrow[r, "f"] & B \arrow[r, "g"] & C \arrow[r] & 0
\end{tikzcd}
\]
qui est exacte en $A$, $B$ et $C$. Cela signifie :
\begin{enumerate}[label=(\roman*)]
\item $f$ est un monomorphisme ($\ker f = 0$) ;
\item $g$ est un épimorphisme ($\coker g = 0$) ;
\item $\im f = \ker g$.
\end{enumerate}
On note souvent une telle suite $0 \to A \to B \to C \to 0$.
\end{definition}

\begin{exemple}\label{ex:ses-Z}
\index{suite exacte!courte!exemples}
\begin{enumerate}[label=(\roman*)]
\item Dans $\Ab$ : $0 \to \Z \xrightarrow{\times n} \Z \to \Z/n\Z \to 0$
  pour tout $n \geq 1$.

\item Dans $R\text{-}\Mod$ : si $N \leq M$ est un sous-module, alors
  $0 \to N \hookrightarrow M \twoheadrightarrow M/N \to 0$ est exacte.

\item La suite $0 \to \Z \xrightarrow{2} \Z \to \Z/2\Z \to 0$
  est une suite exacte courte qui \emph{ne scinde pas}\index{suite exacte!scindée}.
\end{enumerate}
\end{exemple}

\begin{definition}[Suite exacte scindée]\label{def:split-ses}
\index{suite exacte!scindée|textbf}
Une suite exacte courte
$0 \to A \xrightarrow{f} B \xrightarrow{g} C \to 0$
est \textbf{scindée} si l'une des conditions équivalentes suivantes
est satisfaite :
\begin{enumerate}[label=(\roman*)]
\item il existe $s \colon C \to B$ tel que $g \circ s = \id_C$
  (section\index{section} de $g$) ;
\item il existe $r \colon B \to A$ tel que $r \circ f = \id_A$
  (rétraction\index{rétraction} de $f$) ;
\item $B \cong A \oplus C$ via un isomorphisme compatible avec $f$ et $g$.
\end{enumerate}
\end{definition}


%% ============================================================
\section{Foncteurs exacts}\label{sec:exact-functors}
%% ============================================================

\begin{definition}[Foncteur exact à gauche, à droite, exact]\label{def:exact-functor}
\index{foncteur!exact}
\index{foncteur!exact à gauche}
\index{foncteur!exact à droite}
Soit $F \colon \mathcal{A} \to \mathcal{B}$ un foncteur additif entre
catégories abéliennes.
\begin{enumerate}[label=(\roman*)]
\item $F$ est \textbf{exact à gauche}\index{foncteur!exact à gauche|textbf}
  s'il transforme toute suite exacte courte
  $0 \to A \to B \to C \to 0$ en une suite exacte
  $0 \to FA \to FB \to FC$.

\item $F$ est \textbf{exact à droite}\index{foncteur!exact à droite|textbf}
  s'il transforme toute suite exacte courte
  $0 \to A \to B \to C \to 0$ en une suite exacte
  $FA \to FB \to FC \to 0$.

\item $F$ est \textbf{exact}\index{foncteur!exact|textbf}
  s'il est exact à gauche et à droite, c'est-à-dire s'il préserve
  les suites exactes courtes.
\end{enumerate}
\end{definition}

\begin{proposition}\label{prop:hom-left-exact}
\index{Hom@$\Hom$!exactitude à gauche}
Pour tout objet $A$ d'une catégorie abélienne $\mathcal{A}$, les
foncteurs
\[
  \Hom_{\mathcal{A}}(A, -) \colon \mathcal{A} \to \Ab
  \qquad \text{et} \qquad
  \Hom_{\mathcal{A}}(-, A) \colon \mathcal{A}^{\op} \to \Ab
\]
sont exacts à gauche.
\end{proposition}

\begin{proof}
Soit $0 \to B' \xrightarrow{f} B \xrightarrow{g} B'' \to 0$ une suite
exacte courte. On doit montrer que
\[
  0 \to \Hom(A, B') \xrightarrow{f_*} \Hom(A, B) \xrightarrow{g_*} \Hom(A, B'')
\]
est exacte.

\emph{Injectivité de $f_*$.}
Si $f \circ h = 0$ pour $h \colon A \to B'$, alors $h$ se factorise
par $\ker f = 0$, donc $h = 0$.

\emph{Exactitude en $\Hom(A, B)$.}
On a $g_* \circ f_* = (g \circ f)_* = 0$. Réciproquement, si
$g \circ \varphi = 0$ pour $\varphi \colon A \to B$, alors $\varphi$
se factorise par $\ker g = \im f$, donc il existe
$\psi \colon A \to B'$ avec $f \circ \psi = \varphi$.
L'argument pour $\Hom(-, A)$ est dual.
\end{proof}

\begin{exemple}\label{ex:tensor-right-exact}
\index{produit tensoriel!exactitude à droite}
Le foncteur $- \otimes_R N \colon R\text{-}\Mod \to \Ab$ est exact
à droite mais pas exact en général. La mesure de sa non-exactitude
à gauche est le foncteur $\operatorname{Tor}$\index{Tor@$\operatorname{Tor}$},
que nous étudierons au chapitre~\ref{ch:derived}.
\end{exemple}


%% ============================================================
\section{Le lemme du serpent}\label{sec:snake-lemma}
%% ============================================================

\begin{theoreme}[Lemme du serpent]\label{thm:snake-lemma}
\index{lemme!du serpent|textbf}
\index{morphisme connectant}
Soit $\mathcal{A}$ une catégorie abélienne. Étant donné un diagramme
commutatif à lignes exactes
\[
\begin{tikzcd}[column sep=large]
  & A' \arrow[r, "f'"] \arrow[d, "a"]
  & B' \arrow[r, "g'"] \arrow[d, "b"]
  & C' \arrow[r] \arrow[d, "c"]
  & 0 \\
0 \arrow[r]
  & A \arrow[r, "f"]
  & B \arrow[r, "g"]
  & C &
\end{tikzcd}
\]
il existe un morphisme connectant\index{morphisme connectant|textbf}
$\delta \colon \ker c \to \coker a$ et une suite exacte
\[
\begin{tikzcd}[column sep=small]
\ker a \arrow[r] & \ker b \arrow[r] & \ker c
  \arrow[r, "\delta"] & \coker a \arrow[r] & \coker b \arrow[r] & \coker c.
\end{tikzcd}
\]
De plus, si $A' \to B'$ est un monomorphisme (\ie la ligne supérieure
est exacte aussi à gauche), alors $\ker a \to \ker b$ est un
monomorphisme. Dualement, si $B \to C$ est un épimorphisme, alors
$\coker b \to \coker c$ est un épimorphisme.
\end{theoreme}

\begin{proof}
Par le théorème de Freyd--Mitchell\index{théorème!de Freyd--Mitchell},
il suffit de démontrer ce résultat dans $R\text{-}\Mod$, où l'on peut
raisonner sur les éléments. Nous donnons néanmoins une preuve
diagrammatique qui peut être formalisée dans toute catégorie abélienne.

\medskip
\textbf{Étape 1 : Morphismes entre noyaux.}
\index{lemme!du serpent!construction}
Puisque $f \circ a \circ \ker a' = b \circ f' \circ \ker a' = b \circ f' \circ \ker a'$
et le carré commute ($b \circ f' = f \circ a$), la restriction de $b$ à
$\ker a$ donne un morphisme $\ker a \to \ker b$ (car
$b \circ (f' \circ \ker a') = f \circ (a \circ \ker a') = f \circ 0 = 0$,
et $f$ est un monomorphisme dans la ligne inférieure lorsque la suite est
exacte à gauche). Plus précisément, pour tout $x \in \ker a$, on a
$a(x) = 0$, donc $b(f'(x)) = f(a(x)) = 0$, ce qui signifie $f'(x) \in \ker b$.
De même, on construit $\ker b \to \ker c$.

\medskip
\textbf{Étape 2 : Construction du morphisme connectant $\delta$.}
Soit $z \in \ker c$. Comme $g'$ est surjectif, il existe $y \in B'$ tel que
$g'(y) = z$. Considérons $b(y) \in B$. On a
\[
  g(b(y)) = c(g'(y)) = c(z) = 0,
\]
donc $b(y) \in \ker g = \im f$, et il existe $x \in A$ (unique car $f$ est
injectif) tel que $f(x) = b(y)$. On pose
\[
  \delta(z) = \text{classe de } x \text{ dans } \coker a = A / \im a.
\]

\emph{Bonne définition.}
Si $y_1$ est un autre relèvement de $z$, alors $y - y_1 \in \ker g' = \im f'$,
donc $y - y_1 = f'(w)$ pour un $w \in A'$. Alors
$b(y) - b(y_1) = b(f'(w)) = f(a(w))$, d'où $x - x_1 = a(w) \in \im a$,
et les classes coïncident.

\medskip
\textbf{Étape 3 : Exactitude en $\ker c$.}
Soit $z \in \ker c$ dans l'image de $\ker b \to \ker c$, c'est-à-dire
$z = g'(y)$ avec $b(y) = 0$. Alors on peut prendre $x = 0$
(car $f(0) = 0 = b(y)$), donc $\delta(z) = 0$.

Réciproquement, si $\delta(z) = 0$, alors $x = a(w)$ pour un $w \in A'$.
Posons $y_0 = y - f'(w)$. Alors $g'(y_0) = g'(y) = z$ et
$b(y_0) = b(y) - b(f'(w)) = f(x) - f(a(w)) = f(x - x) = 0$,
donc $y_0 \in \ker b$ et $z$ provient de $\ker b$.

\medskip
\textbf{Étape 4 : Exactitude en $\coker a$.}
Soit $\bar{x} \in \coker a$ l'image d'un $x \in A$ par la projection.
Alors $\bar{x} \in \im \delta$ si et seulement si il existe $y \in B'$
avec $f(x) = b(y)$, c'est-à-dire $\bar{x}$ est dans le noyau de
$\coker a \to \coker b$ (car $f(x) = b(y)$ signifie que $f(x) \in \im b$).
D'où l'exactitude en $\coker a$.

Les vérifications de l'exactitude en $\ker a$, $\ker b$, $\coker b$ et
$\coker c$ sont similaires et laissées au lecteur (exercice~\ref{exo:snake-details}).
\end{proof}

Le diagramme complet du lemme du serpent est souvent représenté ainsi :

\[
\begin{tikzcd}[column sep=large, row sep=large]
  & \ker a \arrow[r] \arrow[d, hook]
  & \ker b \arrow[r] \arrow[d, hook]
  & \ker c \arrow[d, hook]
    \arrow[dll, "\delta"',
      rounded corners,
      to path={ -- ([xshift=2.5em]\tikztostart.east)
                |- ([yshift=-1.2em]\tikztostart.south) -| ([xshift=-2.5em]\tikztotarget.west)
                -- (\tikztotarget)}] \\
  & A' \arrow[r, "f'"] \arrow[d, "a"]
  & B' \arrow[r, "g'"] \arrow[d, "b"]
  & C' \arrow[r] \arrow[d, "c"]
  & 0 \\
0 \arrow[r]
  & A \arrow[r, "f"] \arrow[d, two heads]
  & B \arrow[r, "g"] \arrow[d, two heads]
  & C \arrow[d, two heads] \\
  & \coker a \arrow[r]
  & \coker b \arrow[r]
  & \coker c
\end{tikzcd}
\]


%% ============================================================
\section{Le lemme des cinq}\label{sec:five-lemma}
%% ============================================================

\begin{theoreme}[Lemme des cinq]\label{thm:five-lemma}
\index{lemme!des cinq|textbf}
\index{five lemma|see{lemme des cinq}}
Soit un diagramme commutatif dans une catégorie abélienne dont les lignes
sont exactes :
\[
\begin{tikzcd}
A_1 \arrow[r, "d_1"] \arrow[d, "\alpha_1"]
& A_2 \arrow[r, "d_2"] \arrow[d, "\alpha_2"]
& A_3 \arrow[r, "d_3"] \arrow[d, "\alpha_3"]
& A_4 \arrow[r, "d_4"] \arrow[d, "\alpha_4"]
& A_5 \arrow[d, "\alpha_5"] \\
B_1 \arrow[r, "e_1"]
& B_2 \arrow[r, "e_2"]
& B_3 \arrow[r, "e_3"]
& B_4 \arrow[r, "e_4"]
& B_5
\end{tikzcd}
\]
\begin{enumerate}[label=(\roman*)]
\item Si $\alpha_1$ est un épimorphisme et $\alpha_2, \alpha_4$ sont des
  monomorphismes, alors $\alpha_3$ est un monomorphisme.
\item Si $\alpha_5$ est un monomorphisme et $\alpha_2, \alpha_4$ sont des
  épimorphismes, alors $\alpha_3$ est un épimorphisme.
\item En particulier, si $\alpha_1, \alpha_2, \alpha_4, \alpha_5$ sont des
  isomorphismes, alors $\alpha_3$ est un isomorphisme.
\end{enumerate}
\end{theoreme}

\begin{proof}
Par le théorème de Freyd--Mitchell, il suffit de raisonner dans
$R\text{-}\Mod$.

\medskip
\textbf{(i) $\alpha_3$ est un monomorphisme.}
Soit $x \in A_3$ tel que $\alpha_3(x) = 0$. Alors
$e_3(\alpha_3(x)) = 0$, donc $\alpha_4(d_3(x)) = 0$
(par commutativité). Comme $\alpha_4$ est un monomorphisme,
$d_3(x) = 0$. Par exactitude de la ligne supérieure en $A_3$,
il existe $y \in A_2$ tel que $d_2(y) = x$.

Considérons $\alpha_2(y) \in B_2$. On a
$e_2(\alpha_2(y)) = \alpha_3(d_2(y)) = \alpha_3(x) = 0$.
Par exactitude de la ligne inférieure en $B_2$, il existe
$z \in B_1$ tel que $e_1(z) = \alpha_2(y)$.
Comme $\alpha_1$ est un épimorphisme, il existe $w \in A_1$
avec $\alpha_1(w) = z$. Alors
\[
  \alpha_2(d_1(w)) = e_1(\alpha_1(w)) = e_1(z) = \alpha_2(y).
\]
Comme $\alpha_2$ est un monomorphisme, $d_1(w) = y$, donc
$x = d_2(y) = d_2(d_1(w)) = 0$ par exactitude ($d_2 \circ d_1 = 0$
car $\im d_1 \subseteq \ker d_2$). Donc $\alpha_3$ est un monomorphisme.

\medskip
\textbf{(ii) $\alpha_3$ est un épimorphisme.}
Argument dual. Soit $b \in B_3$. On a
$\alpha_4^{-1}(e_3(b))$... plus précisément, $e_3(b) \in B_4$
et comme $\alpha_4$ est un épimorphisme, il existe $a_4 \in A_4$
avec $\alpha_4(a_4) = e_3(b)$.

Considérons $e_4(e_3(b)) = 0$ (par exactitude). Donc
$e_4(\alpha_4(a_4)) = 0$, ce qui donne $\alpha_5(d_4(a_4)) = 0$.
Comme $\alpha_5$ est un monomorphisme, $d_4(a_4) = 0$.
Par exactitude en $A_4$, il existe $a_3 \in A_3$ avec $d_3(a_3) = a_4$.

Posons $b' = b - \alpha_3(a_3)$. Alors
$e_3(b') = e_3(b) - e_3(\alpha_3(a_3)) = e_3(b) - \alpha_4(d_3(a_3))
= e_3(b) - \alpha_4(a_4) = 0$.
Par exactitude en $B_3$, il existe $b_2 \in B_2$ avec $e_2(b_2) = b'$.
Comme $\alpha_2$ est un épimorphisme, il existe $a_2 \in A_2$ avec
$\alpha_2(a_2) = b_2$. Alors
$\alpha_3(d_2(a_2)) = e_2(\alpha_2(a_2)) = e_2(b_2) = b'$.
D'où $b = \alpha_3(a_3) + \alpha_3(d_2(a_2)) = \alpha_3(a_3 + d_2(a_2))$,
et $\alpha_3$ est surjectif.
\end{proof}

\begin{corollaire}[Lemme court des cinq]\label{cor:short-five}
\index{lemme!des cinq!court}
Soit un diagramme commutatif à lignes exactes courtes :
\[
\begin{tikzcd}
0 \arrow[r] & A' \arrow[r] \arrow[d, "\alpha"] & B' \arrow[r] \arrow[d, "\beta"]
  & C' \arrow[r] \arrow[d, "\gamma"] & 0 \\
0 \arrow[r] & A \arrow[r] & B \arrow[r] & C \arrow[r] & 0
\end{tikzcd}
\]
Si $\alpha$ et $\gamma$ sont des isomorphismes, alors $\beta$ est un
isomorphisme.
\end{corollaire}

\begin{proof}
On complète par $0 \to 0$ à gauche et $0 \to 0$ à droite pour obtenir
un diagramme de cinq colonnes, puis on applique le lemme des cinq avec
$\alpha_1 = \alpha_5 = \id_0$, $\alpha_2 = \alpha$, $\alpha_4 = \gamma$.
\end{proof}


%% ============================================================
\section{Objets injectifs et projectifs}\label{sec:injective-projective}
%% ============================================================

\begin{definition}[Objet projectif]\label{def:projective-object}
\index{objet projectif|textbf}
Un objet $P$ d'une catégorie abélienne $\mathcal{A}$ est \textbf{projectif}
si le foncteur $\Hom_{\mathcal{A}}(P, -)$ est exact, c'est-à-dire si
pour tout épimorphisme $g \colon B \twoheadrightarrow C$ et tout morphisme
$f \colon P \to C$, il existe $\tilde{f} \colon P \to B$ tel que
$g \circ \tilde{f} = f$ :
\[
\begin{tikzcd}
& P \arrow[d, "f"] \arrow[dl, dashed, "\tilde{f}"'] \\
B \arrow[r, two heads, "g"] & C \arrow[r] & 0
\end{tikzcd}
\]
\end{definition}

\begin{definition}[Objet injectif]\label{def:injective-object}
\index{objet injectif|textbf}
Dualement, un objet $I$ est \textbf{injectif} si $\Hom_{\mathcal{A}}(-, I)$
est exact, c'est-à-dire si pour tout monomorphisme $f \colon A \hookrightarrow B$
et tout morphisme $g \colon A \to I$, il existe $\tilde{g} \colon B \to I$
avec $\tilde{g} \circ f = g$ :
\[
\begin{tikzcd}
0 \arrow[r] & A \arrow[r, hook, "f"] \arrow[d, "g"'] & B \arrow[dl, dashed, "\tilde{g}"] \\
& I &
\end{tikzcd}
\]
\end{definition}

\begin{proposition}\label{prop:free-projective}
\index{module libre!est projectif}
\index{objet projectif!modules libres}
Tout $R$-module libre est projectif. Plus généralement, les
facteurs directs de modules libres (modules projectifs) sont projectifs.
\end{proposition}

\begin{proof}
Soit $F = \bigoplus_{i \in I} R$ un module libre. Un morphisme
$f \colon F \to C$ est déterminé par les images $f(e_i) \in C$ des
générateurs. Si $g \colon B \twoheadrightarrow C$ est surjectif, on
choisit pour chaque $i$ un $b_i \in B$ avec $g(b_i) = f(e_i)$.
La propriété universelle du module libre donne un unique
$\tilde{f} \colon F \to B$ avec $\tilde{f}(e_i) = b_i$, et
$g \circ \tilde{f} = f$.
\end{proof}

\begin{theoreme}[Critère de Baer]\label{thm:baer-criterion}
\index{critère de Baer|textbf}
\index{Baer, Reinhold}
Un $R$-module $I$ est injectif si et seulement si pour tout
idéal $J \subseteq R$ et tout morphisme $f \colon J \to I$,
il existe une extension $\tilde{f} \colon R \to I$ avec
$\tilde{f}|_J = f$.
\end{theoreme}

\begin{proof}
La nécessité est immédiate (prendre le monomorphisme $J \hookrightarrow R$).
Pour la suffisance, soit $f \colon A \hookrightarrow B$ un monomorphisme
et $g \colon A \to I$ un morphisme. On utilise le lemme de Zorn sur
l'ensemble des extensions partielles de $g$ le long de sous-modules
de $B$ contenant $A$. Un élément maximal $(M, h)$ doit avoir $M = B$ :
sinon, on prend $b \in B \setminus M$ et on pose
$J = \{r \in R : rb \in M\}$. La composée
$J \xrightarrow{r \mapsto rb} M \xrightarrow{h} I$
s'étend à $R$ par hypothèse, ce qui permet d'étendre $h$ à $M + Rb$,
contredisant la maximalité.
\end{proof}

\begin{definition}[Assez d'injectifs / de projectifs]\label{def:enough-injectives}
\index{assez d'injectifs|textbf}
\index{assez de projectifs|textbf}
Une catégorie abélienne $\mathcal{A}$ a \textbf{assez d'injectifs}
si tout objet $A$ admet un monomorphisme $A \hookrightarrow I$ avec
$I$ injectif.

Dualement, $\mathcal{A}$ a \textbf{assez de projectifs} si tout
objet $A$ admet un épimorphisme $P \twoheadrightarrow A$ avec $P$
projectif.
\end{definition}

\begin{theoreme}\label{thm:R-mod-enough}
\index{R-Mod@$R\text{-}\Mod$!assez d'injectifs}
\index{R-Mod@$R\text{-}\Mod$!assez de projectifs}
Pour tout anneau $R$, la catégorie $R\text{-}\Mod$ a assez d'injectifs
et assez de projectifs.
\end{theoreme}

\begin{proof}
\emph{Assez de projectifs.} Tout module $M$ est quotient d'un module
libre (prendre le module libre sur un ensemble de générateurs de $M$),
et les modules libres sont projectifs.

\emph{Assez d'injectifs.} On utilise le plongement de Baer. Tout
groupe abélien $A$ se plonge dans un groupe abélien divisible\index{groupe abélien!divisible}
(donc injectif en tant que $\Z$-module). Pour un $R$-module $M$,
on plonge le $\Z$-module sous-jacent dans un $\Z$-module injectif $Q$,
puis on montre que $\Hom_\Z(R, Q)$ est un $R$-module injectif
et que l'application $M \to \Hom_\Z(R, Q)$ définie par
$m \mapsto (r \mapsto rm \in Q)$ est un monomorphisme de $R$-modules.
\end{proof}


%% ============================================================
\section{Résolutions}\label{sec:resolutions}
%% ============================================================

\begin{definition}[Résolution projective]\label{def:projective-resolution}
\index{résolution!projective|textbf}
Soit $A$ un objet d'une catégorie abélienne $\mathcal{A}$.
Une \textbf{résolution projective} de $A$ est une suite exacte
\[
\begin{tikzcd}
\cdots \arrow[r] & P_2 \arrow[r, "d_2"] & P_1 \arrow[r, "d_1"]
  & P_0 \arrow[r, "\varepsilon"] & A \arrow[r] & 0
\end{tikzcd}
\]
où chaque $P_n$ est projectif. On note $P_\bullet \to A$ une telle résolution.
\end{definition}

\begin{definition}[Résolution injective]\label{def:injective-resolution}
\index{résolution!injective|textbf}
Dualement, une \textbf{résolution injective} de $A$ est une suite exacte
\[
\begin{tikzcd}
0 \arrow[r] & A \arrow[r, "\eta"] & I^0 \arrow[r, "d^0"]
  & I^1 \arrow[r, "d^1"] & I^2 \arrow[r] & \cdots
\end{tikzcd}
\]
où chaque $I^n$ est injectif. On note $A \to I^\bullet$.
\end{definition}

\begin{proposition}\label{prop:resolution-existence}
\index{résolution!existence}
Si $\mathcal{A}$ a assez de projectifs (resp.\ d'injectifs), tout
objet admet une résolution projective (resp.\ injective).
\end{proposition}

\begin{proof}
Pour une résolution projective : on choisit un épimorphisme
$P_0 \twoheadrightarrow A$ avec $P_0$ projectif, puis un épimorphisme
$P_1 \twoheadrightarrow \ker(P_0 \to A)$ avec $P_1$ projectif, et
ainsi de suite par récurrence. L'argument pour les résolutions
injectives est dual.
\end{proof}

\begin{theoreme}[Lemme de comparaison]\label{thm:comparison-lemma}
\index{lemme!de comparaison|textbf}
Soit $f \colon A \to B$ un morphisme, $P_\bullet \to A$ une résolution
projective et $Q_\bullet \to B$ une suite exacte quelconque (pas
nécessairement projective). Alors il existe un morphisme de chaînes
$\varphi_\bullet \colon P_\bullet \to Q_\bullet$ relevant $f$, et
deux tels relèvements sont homotopes\index{homotopie de chaînes}.
\end{theoreme}

\begin{proof}
On construit $\varphi_n \colon P_n \to Q_n$ par récurrence.
Pour $n = 0$ : on a $\varepsilon_B \circ \varphi_0 = f \circ \varepsilon_A$,
et comme $P_0$ est projectif et $\varepsilon_B$ est surjectif, le
relèvement existe. Pour le pas de récurrence : on utilise la
projectivité de $P_n$ et l'exactitude de $Q_\bullet$ pour construire
$\varphi_n$ relevant le morphisme induit sur les noyaux.

L'unicité à homotopie près se démontre de manière similaire, en
construisant une homotopie $s_n \colon P_n \to Q_{n+1}$ par
récurrence, utilisant à chaque étape la projectivité de $P_n$.
\end{proof}


%% ============================================================
\section{Exercices}\label{sec:exos-abelian}
%% ============================================================

\begin{exercice}\label{exo:additive-functor-zero}
\index{foncteur!additif!préserve l'objet nul}
Montrer qu'un foncteur additif $F \colon \mathcal{A} \to \mathcal{B}$
entre catégories additives préserve l'objet nul et les biproduits.
\end{exercice}

\begin{exercice}\label{exo:abelian-mono-epi}
\index{monomorphisme!dans une catégorie abélienne}
Soit $\mathcal{A}$ une catégorie abélienne et $f \colon A \to B$ un
morphisme. Montrer que :
\begin{enumerate}[label=(\alph*)]
\item $f$ est un monomorphisme si et seulement si $\ker f = 0$ ;
\item $f$ est un épimorphisme si et seulement si $\coker f = 0$ ;
\item $f$ est un isomorphisme si et seulement si $\ker f = 0$ et
  $\coker f = 0$.
\end{enumerate}
\end{exercice}

\begin{exercice}\label{exo:snake-details}
\index{lemme!du serpent!détails}
Compléter les détails de la preuve du lemme du serpent
(théorème~\ref{thm:snake-lemma}) en vérifiant l'exactitude de la suite
$\ker a \to \ker b \to \ker c$ et $\coker a \to \coker b \to \coker c$.
\end{exercice}

\begin{exercice}\label{exo:horseshoe}
\index{lemme!du fer à cheval}
\textbf{(Lemme du fer à cheval.)}
Soit $0 \to A' \to A \to A'' \to 0$ une suite exacte courte dans
une catégorie abélienne ayant assez de projectifs. Si $P'_\bullet \to A'$
et $P''_\bullet \to A''$ sont des résolutions projectives, montrer qu'il
existe une résolution projective $P_\bullet \to A$ avec
$P_n = P'_n \oplus P''_n$ qui s'insère dans une suite exacte courte
de complexes $0 \to P'_\bullet \to P_\bullet \to P''_\bullet \to 0$.
\end{exercice}

\begin{exercice}\label{exo:ses-projective}
\index{suite exacte!courte!scindée!caractérisation projective}
Soit $0 \to A \xrightarrow{f} B \xrightarrow{g} C \to 0$ une suite
exacte courte dans $R\text{-}\Mod$. Montrer que les conditions
suivantes sont équivalentes :
\begin{enumerate}[label=(\roman*)]
\item la suite est scindée ;
\item $g$ admet une section ;
\item $f$ admet une rétraction ;
\item pour tout module $M$, la suite
  $0 \to \Hom(M, A) \to \Hom(M, B) \to \Hom(M, C) \to 0$ est exacte.
\end{enumerate}
En déduire qu'un module $C$ est projectif si et seulement si toute suite
exacte courte $0 \to A \to B \to C \to 0$ est scindée.
\end{exercice}

\begin{exercice}\label{exo:injective-divisible}
\index{groupe abélien!divisible}
\index{objet injectif!groupes abéliens divisibles}
Un groupe abélien $A$ est dit \textbf{divisible}\index{divisible} si pour
tout $n \geq 1$ et tout $a \in A$, il existe $b \in A$ avec $nb = a$.
\begin{enumerate}[label=(\alph*)]
\item Montrer que $\Q$ et $\Q/\Z$ sont divisibles.
\item Montrer qu'un groupe abélien est injectif (dans $\Ab$) si et seulement
  s'il est divisible.
  \textit{Indication : utiliser le critère de Baer.}
\item En déduire que $\Ab$ a assez d'injectifs.
\end{enumerate}
\end{exercice}

\begin{exercice}\label{exo:nine-lemma}
\index{lemme!des neuf}
\textbf{(Lemme $3 \times 3$.)}
Soit un diagramme commutatif dans une catégorie abélienne dont les
colonnes et les deux premières lignes sont exactes :
\[
\begin{tikzcd}
& 0 \arrow[d] & 0 \arrow[d] & 0 \arrow[d] \\
0 \arrow[r] & A' \arrow[r] \arrow[d] & B' \arrow[r] \arrow[d] & C' \arrow[r] \arrow[d] & 0 \\
0 \arrow[r] & A \arrow[r] \arrow[d] & B \arrow[r] \arrow[d] & C \arrow[r] \arrow[d] & 0 \\
0 \arrow[r] & A'' \arrow[r] \arrow[d] & B'' \arrow[r] \arrow[d] & C'' \arrow[r] \arrow[d] & 0 \\
& 0 & 0 & 0
\end{tikzcd}
\]
Montrer que la troisième ligne est également exacte.
\textit{Indication : appliquer deux fois le lemme du serpent.}
\end{exercice}

\begin{exercice}\label{exo:grothendieck-AB5}
\index{axiome AB5}
\index{catégorie!de Grothendieck}
Soit $\mathcal{A}$ une catégorie abélienne. On dit que $\mathcal{A}$
satisfait l'axiome \textbf{AB5}\index{AB5} si $\mathcal{A}$ admet
toutes les colimites filtrantes et si les colimites filtrantes
commutent aux suites exactes. Montrer que $R\text{-}\Mod$ satisfait AB5.
\end{exercice}

\index{suite exacte|)}
\index{catégorie!abélienne|)}


%%%%%%%%%%%%%%%%%%%%%%%%%%%%%%%%%%%%%%%%%%%%%%%%%%%%%%%%%%%%
%%           CHAPITRE 8 : CATÉGORIES DÉRIVÉES             %%
%%                   — INTRODUCTION                        %%
%%%%%%%%%%%%%%%%%%%%%%%%%%%%%%%%%%%%%%%%%%%%%%%%%%%%%%%%%%%%

\chapter{Catégories dérivées --- Introduction}\label{ch:derived}
\minitoc

\vspace{1em}

Les catégories dérivées\index{catégorie!dérivée|(}, introduites par
Grothendieck\index{Grothendieck, Alexander} et développées par
Verdier\index{Verdier, Jean-Louis} dans sa thèse (1967), fournissent le
cadre conceptuel adéquat pour l'algèbre homologique. L'idée fondatrice est
simple mais profonde : au lieu de travailler avec un objet $A$ et ses
résolutions, on travaille directement avec la catégorie des complexes de
chaînes, en « inversant » les quasi-isomorphismes. Ce processus de
\emph{localisation}\index{localisation!de catégories} produit la catégorie
dérivée $D(\mathcal{A})$, dans laquelle un objet et sa résolution deviennent
isomorphes. La structure supplémentaire de \emph{catégorie
triangulée}\index{catégorie!triangulée} encode l'information des suites
exactes longues.

\index{complexe de chaînes|(}

%% ============================================================
\section{Complexes de chaînes}\label{sec:chain-complexes}
%% ============================================================

\begin{definition}[Complexe de chaînes]\label{def:chain-complex}
\index{complexe de chaînes|textbf}
\index{différentielle}
Soit $\mathcal{A}$ une catégorie abélienne. Un \textbf{complexe de chaînes}
(ou \textbf{complexe cochain}) dans $\mathcal{A}$ est une suite
$(A^n, d^n)_{n \in \Z}$ d'objets et de morphismes
\[
\begin{tikzcd}
\cdots \arrow[r] & A^{n-1} \arrow[r, "d^{n-1}"] & A^n \arrow[r, "d^n"]
  & A^{n+1} \arrow[r] & \cdots
\end{tikzcd}
\]
telle que $d^{n} \circ d^{n-1} = 0$ pour tout $n \in \Z$.
Les morphismes $d^n$ sont appelés les \textbf{différentielles}\index{différentielle|textbf}.
On note le complexe $A^\bullet$ ou $(A^\bullet, d)$.
\end{definition}

\begin{definition}[Morphisme de complexes]\label{def:chain-map}
\index{morphisme!de complexes|textbf}
Un \textbf{morphisme de complexes} $f \colon A^\bullet \to B^\bullet$
est une famille de morphismes $(f^n \colon A^n \to B^n)_{n \in \Z}$
telle que le diagramme suivant commute pour tout $n$ :
\[
\begin{tikzcd}
A^n \arrow[r, "d_A^n"] \arrow[d, "f^n"'] & A^{n+1} \arrow[d, "f^{n+1}"] \\
B^n \arrow[r, "d_B^n"] & B^{n+1}
\end{tikzcd}
\]
c'est-à-dire $f^{n+1} \circ d_A^n = d_B^n \circ f^n$.
\end{definition}

\begin{notation}\label{not:Ch}
\index{Ch(A)@$\mathsf{Ch}(\mathcal{A})$}
On note $\mathsf{Ch}(\mathcal{A})$ la catégorie dont les objets sont les
complexes de chaînes dans $\mathcal{A}$ et les morphismes sont les
morphismes de complexes. On note aussi :
\begin{itemize}
\item $\mathsf{Ch}^+(\mathcal{A})$\index{Ch+@$\mathsf{Ch}^+(\mathcal{A})$} :
  complexes bornés inférieurement ($A^n = 0$ pour $n \ll 0$) ;
\item $\mathsf{Ch}^-(\mathcal{A})$\index{Ch-@$\mathsf{Ch}^-(\mathcal{A})$} :
  complexes bornés supérieurement ($A^n = 0$ pour $n \gg 0$) ;
\item $\mathsf{Ch}^b(\mathcal{A})$\index{Chb@$\mathsf{Ch}^b(\mathcal{A})$} :
  complexes bornés ($A^n = 0$ pour $|n| \gg 0$).
\end{itemize}
\end{notation}

\begin{definition}[Cohomologie d'un complexe]\label{def:cohomology}
\index{cohomologie!d'un complexe|textbf}
\index{cocycle}
\index{cobord}
Le $n$-ième \textbf{objet de cohomologie}\index{cohomologie|textbf}
du complexe $A^\bullet$ est
\[
  H^n(A^\bullet) = \ker d^n / \im d^{n-1}
  = \ker(d^n \colon A^n \to A^{n+1}) \big/ \im(d^{n-1} \colon A^{n-1} \to A^n).
\]
Ceci a un sens car $\im d^{n-1} \subseteq \ker d^n$ (puisque $d^n \circ d^{n-1} = 0$).
Les éléments de $\ker d^n$ sont les \textbf{$n$-cocycles}\index{cocycle|textbf}
et ceux de $\im d^{n-1}$ les \textbf{$n$-cobords}\index{cobord|textbf}.
\end{definition}

\begin{proposition}\label{prop:cohomology-functor}
\index{cohomologie!foncteur}
La cohomologie définit un foncteur
$H^n \colon \mathsf{Ch}(\mathcal{A}) \to \mathcal{A}$
pour chaque $n \in \Z$.
\end{proposition}

\begin{proof}
Soit $f \colon A^\bullet \to B^\bullet$ un morphisme de complexes.
Comme $f^n$ envoie $\ker d_A^n$ dans $\ker d_B^n$ (par commutativité)
et $\im d_A^{n-1}$ dans $\im d_B^{n-1}$, il induit un morphisme
$H^n(f) \colon H^n(A^\bullet) \to H^n(B^\bullet)$. La fonctorialité
($H^n(\id) = \id$ et $H^n(g \circ f) = H^n(g) \circ H^n(f)$) est
immédiate.
\end{proof}

\begin{proposition}\label{prop:Ch-abelian}
\index{Ch(A)@$\mathsf{Ch}(\mathcal{A})$!catégorie abélienne}
Si $\mathcal{A}$ est abélienne, alors $\mathsf{Ch}(\mathcal{A})$ est
abélienne, les noyaux, conoyaux et (co)images étant calculés degré
par degré.
\end{proposition}

\begin{proof}
On vérifie les axiomes. L'objet nul est le complexe nul.
Le biproduit est $(A \oplus B)^n = A^n \oplus B^n$ avec
$d_{A \oplus B}^n = d_A^n \oplus d_B^n$.
Le noyau de $f \colon A^\bullet \to B^\bullet$ est le complexe
$(\ker f^n, d^n|_{\ker f^n})_{n \in \Z}$, qui est bien un sous-complexe
car $d^n$ envoie $\ker f^n$ dans $\ker f^{n+1}$. De même pour le
conoyau. La condition image $=$ coimage se vérifie degré par degré
car elle est satisfaite dans $\mathcal{A}$.
\end{proof}


%% ============================================================
\section{Homotopie de chaînes et catégorie homotopique}%
\label{sec:chain-homotopy}
%% ============================================================

\begin{definition}[Homotopie de chaînes]\label{def:chain-homotopy}
\index{homotopie de chaînes|textbf}
Soient $f, g \colon A^\bullet \to B^\bullet$ deux morphismes de complexes.
Une \textbf{homotopie de chaînes} de $f$ vers $g$ est une famille de
morphismes $(s^n \colon A^n \to B^{n-1})_{n \in \Z}$ telle que
\[
  f^n - g^n = d_B^{n-1} \circ s^n + s^{n+1} \circ d_A^n
  \qquad \text{pour tout } n \in \Z.
\]
\[
\begin{tikzcd}
\cdots \arrow[r]
& A^{n-1} \arrow[r, "d_A^{n-1}"] \arrow[d]
& A^n \arrow[r, "d_A^n"] \arrow[d] \arrow[dl, dashed, "s^n"']
& A^{n+1} \arrow[r] \arrow[d] \arrow[dl, dashed, "s^{n+1}"']
& \cdots \\
\cdots \arrow[r]
& B^{n-1} \arrow[r, "d_B^{n-1}"]
& B^n \arrow[r, "d_B^n"]
& B^{n+1} \arrow[r]
& \cdots
\end{tikzcd}
\]
On dit que $f$ et $g$ sont \textbf{homotopes}\index{homotope} et on
note $f \sim g$.
\end{definition}

\begin{proposition}\label{prop:homotopy-equivalence}
\index{homotopie de chaînes!relation d'équivalence}
La relation d'homotopie $\sim$ est une relation d'équivalence sur
$\Hom_{\mathsf{Ch}(\mathcal{A})}(A^\bullet, B^\bullet)$, compatible
avec la composition.
\end{proposition}

\begin{proof}
\emph{Réflexivité :} $f \sim f$ avec $s^n = 0$.
\emph{Symétrie :} si $s$ est une homotopie de $f$ vers $g$, alors $-s$ est
une homotopie de $g$ vers $f$.
\emph{Transitivité :} si $s$ est une homotopie de $f$ vers $g$ et $t$ de
$g$ vers $h$, alors $s + t$ est une homotopie de $f$ vers $h$.

\emph{Compatibilité.} Si $f \sim g \colon A^\bullet \to B^\bullet$
via $s$ et $h \colon B^\bullet \to C^\bullet$, alors $h \circ f \sim h \circ g$
via $h \circ s$. Si $h \colon C^\bullet \to A^\bullet$, alors
$f \circ h \sim g \circ h$ via $s \circ h$.
\end{proof}

\begin{proposition}\label{prop:homotopic-maps-cohomology}
\index{homotopie de chaînes!et cohomologie}
Si $f \sim g \colon A^\bullet \to B^\bullet$, alors
$H^n(f) = H^n(g)$ pour tout $n$.
\end{proposition}

\begin{proof}
Soit $[a] \in H^n(A^\bullet)$, c'est-à-dire $a \in \ker d_A^n$.
Alors
\[
  f^n(a) - g^n(a) = d_B^{n-1}(s^n(a)) + s^{n+1}(d_A^n(a))
  = d_B^{n-1}(s^n(a)) + 0 \in \im d_B^{n-1},
\]
donc $[f^n(a)] = [g^n(a)]$ dans $H^n(B^\bullet)$.
\end{proof}

\begin{definition}[Catégorie homotopique]\label{def:homotopy-category}
\index{catégorie!homotopique|textbf}
\index{K(A)@$K(\mathcal{A})$}
La \textbf{catégorie homotopique} $K(\mathcal{A})$ est la catégorie dont :
\begin{itemize}
\item les objets sont les complexes de chaînes dans $\mathcal{A}$ ;
\item les morphismes sont les classes d'homotopie :
  $\Hom_{K(\mathcal{A})}(A^\bullet, B^\bullet)
  = \Hom_{\mathsf{Ch}(\mathcal{A})}(A^\bullet, B^\bullet) / {\sim}$.
\end{itemize}
On définit de même $K^+(\mathcal{A})$, $K^-(\mathcal{A})$,
$K^b(\mathcal{A})$\index{K+@$K^+(\mathcal{A})$}%
\index{K-@$K^-(\mathcal{A})$}\index{Kb@$K^b(\mathcal{A})$}.
\end{definition}

\begin{definition}[Équivalence d'homotopie]\label{def:homotopy-equivalence}
\index{équivalence d'homotopie|textbf}
Un morphisme de complexes $f \colon A^\bullet \to B^\bullet$ est une
\textbf{équivalence d'homotopie} s'il devient un isomorphisme dans
$K(\mathcal{A})$, c'est-à-dire s'il existe $g \colon B^\bullet \to A^\bullet$
avec $g \circ f \sim \id_{A^\bullet}$ et $f \circ g \sim \id_{B^\bullet}$.
\end{definition}


%% ============================================================
\section{Quasi-isomorphismes et localisation}\label{sec:quasi-iso}
%% ============================================================

\begin{definition}[Quasi-isomorphisme]\label{def:quasi-iso}
\index{quasi-isomorphisme|textbf}
Un morphisme de complexes $f \colon A^\bullet \to B^\bullet$ est un
\textbf{quasi-isomorphisme} si
$H^n(f) \colon H^n(A^\bullet) \xrightarrow{\;\sim\;} H^n(B^\bullet)$
est un isomorphisme pour tout $n \in \Z$.
\end{definition}

\begin{remarque}\label{rem:htpy-eq-is-qiso}
\index{équivalence d'homotopie!est un quasi-isomorphisme}
Toute équivalence d'homotopie est un quasi-isomorphisme (par la
proposition~\ref{prop:homotopic-maps-cohomology}), mais la réciproque
est fausse en général.
\end{remarque}

\begin{exemple}\label{ex:quasi-iso-not-htpy}
\index{quasi-isomorphisme!qui n'est pas une équivalence d'homotopie}
Dans $\Ab$, considérons les complexes
\[
  A^\bullet : \quad 0 \to \Z \xrightarrow{2} \Z \to 0
  \qquad \text{et} \qquad
  B^\bullet : \quad 0 \to 0 \to \Z/2\Z \to 0
\]
(concentrés en degrés 0 et 1). Le morphisme $f \colon A^\bullet \to B^\bullet$
donné par $f^0 = 0$ et $f^1 = \text{projection}$ est un quasi-isomorphisme
($H^0 = 0$, $H^1 = \Z/2\Z$ pour les deux), mais \emph{pas} une
équivalence d'homotopie.
\end{exemple}

\begin{definition}[Localisation d'une catégorie]\label{def:localization}
\index{localisation!de catégories|textbf}
\index{Gabriel, Pierre}
\index{Zisman, Michel}
Soit $\mathcal{C}$ une catégorie et $S$ une classe de morphismes.
La \textbf{localisation} $\mathcal{C}[S^{-1}]$ est une catégorie munie
d'un foncteur $Q \colon \mathcal{C} \to \mathcal{C}[S^{-1}]$ tel que :
\begin{enumerate}[label=(\roman*)]
\item $Q(s)$ est un isomorphisme pour tout $s \in S$ ;
\item $Q$ est universel pour cette propriété : pour tout foncteur
  $F \colon \mathcal{C} \to \mathcal{D}$ tel que $F(s)$ est un
  isomorphisme pour tout $s \in S$, il existe un unique foncteur
  $\bar{F} \colon \mathcal{C}[S^{-1}] \to \mathcal{D}$ avec
  $\bar{F} \circ Q = F$.
\end{enumerate}
\end{definition}

\begin{remarque}\label{rem:localization-existence}
\index{localisation!existence}
La localisation existe toujours (modulo des problèmes ensemblistes),
mais les morphismes dans $\mathcal{C}[S^{-1}]$ sont en général
des « fractions » formelles\index{fraction!formelle}
$s_1^{-1} \circ f_1 \circ s_2^{-1} \circ f_2 \circ \cdots$,
ce qui rend la catégorie difficile à manipuler. La notion de
\emph{système multiplicatif}\index{système multiplicatif} permet de
simplifier cette description.
\end{remarque}

\begin{definition}[Système multiplicatif]\label{def:multiplicative-system}
\index{système multiplicatif|textbf}
\index{conditions d'Ore}
Une classe $S$ de morphismes dans une catégorie $\mathcal{C}$ est un
\textbf{système multiplicatif} si :
\begin{enumerate}[label=(MS\arabic*)]
\item\label{MS1} $S$ contient toutes les identités et est stable par
  composition ;
\item\label{MS2} \textbf{(Condition d'Ore à gauche et à droite.)}
  Pour tout diagramme $A \xrightarrow{f} C \xleftarrow{s} B$ avec
  $s \in S$, il existe un diagramme $A \xleftarrow{t} D \xrightarrow{g} B$
  avec $t \in S$ et $f \circ t = s \circ g$ (et dualement) :
\[
\begin{tikzcd}
D \arrow[r, "g"] \arrow[d, "t"', "{\in S}"] & B \arrow[d, "s", "{\in S}"'] \\
A \arrow[r, "f"] & C
\end{tikzcd}
\]
\item\label{MS3} Pour $f, g \colon A \to B$, il existe $s \in S$ avec
  $s \circ f = s \circ g$ si et seulement s'il existe $t \in S$ avec
  $f \circ t = g \circ t$.
\end{enumerate}
\end{definition}

\begin{proposition}\label{prop:qiso-multiplicative}
\index{quasi-isomorphisme!système multiplicatif}
La classe des quasi-isomorphismes dans $K(\mathcal{A})$ forme un
système multiplicatif.
\end{proposition}

\begin{proof}
La vérification de \ref{MS1} est immédiate.
Pour \ref{MS2} et \ref{MS3}, on utilise la construction du
\emph{cône de mapping}\index{cône!de mapping} (voir ci-dessous) et
les propriétés de la catégorie homotopique. Nous omettons les détails
techniques.
\end{proof}


%% ============================================================
\section{La catégorie dérivée}\label{sec:derived-category}
%% ============================================================

\begin{definition}[Catégorie dérivée]\label{def:derived-category}
\index{catégorie!dérivée|textbf}
\index{D(A)@$D(\mathcal{A})$}
La \textbf{catégorie dérivée} $D(\mathcal{A})$ est la localisation de la
catégorie homotopique $K(\mathcal{A})$ par rapport aux quasi-isomorphismes :
\[
  D(\mathcal{A}) = K(\mathcal{A})[\mathsf{Qis}^{-1}].
\]
On définit de même les variantes bornées :
\begin{itemize}
\item $D^+(\mathcal{A}) = K^+(\mathcal{A})[\mathsf{Qis}^{-1}]$\index{D+@$D^+(\mathcal{A})$} ;
\item $D^-(\mathcal{A}) = K^-(\mathcal{A})[\mathsf{Qis}^{-1}]$\index{D-@$D^-(\mathcal{A})$} ;
\item $D^b(\mathcal{A}) = K^b(\mathcal{A})[\mathsf{Qis}^{-1}]$\index{Db@$D^b(\mathcal{A})$}.
\end{itemize}
\end{definition}

\begin{remarque}\label{rem:derived-category-morphisms}
\index{catégorie!dérivée!morphismes}
\index{toit}
Un morphisme $A^\bullet \to B^\bullet$ dans $D(\mathcal{A})$ est
représenté par un \textbf{toit}\index{toit|textbf} :
\[
\begin{tikzcd}
& C^\bullet \arrow[dl, "s"', "\sim"] \arrow[dr, "f"] \\
A^\bullet & & B^\bullet
\end{tikzcd}
\]
où $s$ est un quasi-isomorphisme. La composition de deux toits
nécessite la condition d'Ore (MS2).
\end{remarque}

\begin{theoreme}\label{thm:derived-via-injectives}
\index{catégorie!dérivée!et résolutions injectives}
Si $\mathcal{A}$ a assez d'injectifs, le foncteur naturel
\[
  K^+(\mathcal{I}) \longrightarrow D^+(\mathcal{A})
\]
est une équivalence de catégories, où $\mathcal{I}$ désigne la
sous-catégorie pleine des objets injectifs de $\mathcal{A}$.
Autrement dit, tout objet de $D^+(\mathcal{A})$ est isomorphe à
un complexe d'injectifs.
\end{theoreme}

\begin{proof}
\emph{Essentielle surjectivité.} Soit $A^\bullet \in D^+(\mathcal{A})$.
Par récurrence, on construit une résolution injective
$I^\bullet$ du complexe $A^\bullet$ : c'est un complexe d'injectifs
muni d'un quasi-isomorphisme $A^\bullet \to I^\bullet$. Cette
construction généralise les résolutions injectives d'un seul objet
(proposition~\ref{prop:resolution-existence}).

\emph{Pleine fidélité.} Il faut montrer que si $I^\bullet, J^\bullet$
sont des complexes bornés inférieurement d'injectifs, alors
\[
  \Hom_{K^+(\mathcal{A})}(I^\bullet, J^\bullet)
  \xrightarrow{\;\sim\;} \Hom_{D^+(\mathcal{A})}(I^\bullet, J^\bullet).
\]
L'injectivité vient du fait qu'un morphisme homotope à zéro vers un
complexe d'injectifs est nul dans la catégorie dérivée. La surjectivité
utilise le fait que tout toit peut être « redressé » en utilisant
l'injectivité des composantes.
\end{proof}


%% ============================================================
\section{Le foncteur de décalage et le cône de mapping}%
\label{sec:shift-cone}
%% ============================================================

\begin{definition}[Foncteur de décalage]\label{def:shift-functor}
\index{foncteur!de décalage|textbf}
\index{décalage}
\index{translation}
Le \textbf{foncteur de décalage} (ou \textbf{translation})
$[1] \colon \mathsf{Ch}(\mathcal{A}) \to \mathsf{Ch}(\mathcal{A})$
est défini par
\[
  (A[1])^n = A^{n+1}, \qquad d_{A[1]}^n = -d_A^{n+1}.
\]
Pour un morphisme $f \colon A^\bullet \to B^\bullet$, on pose
$(f[1])^n = f^{n+1}$. Plus généralement, $A[k]^n = A^{n+k}$ avec
$d_{A[k]}^n = (-1)^k d_A^{n+k}$.
\end{definition}

\begin{definition}[Cône de mapping]\label{def:mapping-cone}
\index{cône!de mapping|textbf}
\index{mapping cone|see{cône de mapping}}
Soit $f \colon A^\bullet \to B^\bullet$ un morphisme de complexes.
Le \textbf{cône de mapping} de $f$, noté $\operatorname{Cone}(f)$,
est le complexe défini par
\[
  \operatorname{Cone}(f)^n = A^{n+1} \oplus B^n
\]
avec la différentielle
\[
  d_{\operatorname{Cone}(f)}^n =
  \begin{pmatrix} -d_A^{n+1} & 0 \\ f^{n+1} & d_B^n \end{pmatrix}
  \colon A^{n+1} \oplus B^n \longrightarrow A^{n+2} \oplus B^{n+1}.
\]
\end{definition}

\begin{proposition}\label{prop:cone-exact-triangle}
\index{cône!de mapping!suite exacte}
\index{triangle!distingué}
Pour tout morphisme de complexes $f \colon A^\bullet \to B^\bullet$,
on a une suite de morphismes de complexes
\[
\begin{tikzcd}
A^\bullet \arrow[r, "f"] & B^\bullet \arrow[r, "\iota"]
  & \operatorname{Cone}(f) \arrow[r, "\pi"] & A[1]^\bullet
\end{tikzcd}
\]
où $\iota^n(b) = (0, b)$ et $\pi^n(a, b) = a$. De plus, cette suite
induit une suite exacte longue en cohomologie :
\[
\begin{tikzcd}[column sep=small]
\cdots \arrow[r]
& H^n(A^\bullet) \arrow[r, "H^n(f)"]
& H^n(B^\bullet) \arrow[r]
& H^n(\operatorname{Cone}(f)) \arrow[r]
& H^{n+1}(A^\bullet) \arrow[r]
& \cdots
\end{tikzcd}
\]
\end{proposition}

\begin{proof}
On vérifie d'abord que $d_{\operatorname{Cone}(f)}^2 = 0$ :
\[
  \begin{pmatrix} -d_A^{n+2} & 0 \\ f^{n+2} & d_B^{n+1} \end{pmatrix}
  \begin{pmatrix} -d_A^{n+1} & 0 \\ f^{n+1} & d_B^n \end{pmatrix}
  = \begin{pmatrix} d_A^{n+2} d_A^{n+1} & 0 \\
    -f^{n+2} d_A^{n+1} + d_B^{n+1} f^{n+1} & d_B^{n+1} d_B^n \end{pmatrix}
  = 0,
\]
utilisant $d^2 = 0$ et $f \circ d_A = d_B \circ f$.

Pour la suite exacte courte de complexes, on a
$0 \to B^\bullet \xrightarrow{\iota} \operatorname{Cone}(f) \xrightarrow{\pi} A[1]^\bullet \to 0$
qui est exacte degré par degré. La suite exacte longue en cohomologie
en découle (voir théorème~\ref{thm:long-exact-sequence}).
Le morphisme connectant $H^n(A[1]^\bullet) = H^{n+1}(A^\bullet) \to H^{n+1}(B^\bullet)$
s'identifie à $H^{n+1}(f)$.
\end{proof}


%% ============================================================
\section{Structure triangulée}\label{sec:triangulated}
%% ============================================================

\begin{definition}[Catégorie triangulée]\label{def:triangulated}
\index{catégorie!triangulée|textbf}
\index{axiomes TR}
Une \textbf{catégorie triangulée} est un triplet
$(\mathcal{T}, [1], \Delta)$ où :
\begin{itemize}
\item $\mathcal{T}$ est une catégorie additive ;
\item $[1] \colon \mathcal{T} \to \mathcal{T}$ est un automorphisme
  (le \textbf{foncteur de translation}\index{foncteur!de translation}) ;
\item $\Delta$ est une classe de suites
  $A \xrightarrow{u} B \xrightarrow{v} C \xrightarrow{w} A[1]$
  appelées \textbf{triangles distingués}\index{triangle!distingué|textbf},
\end{itemize}
satisfaisant les axiomes suivants.
\end{definition}

\begin{definition}[Axiomes TR1--TR4]\label{def:TR-axioms}
\index{TR1}\index{TR2}\index{TR3}\index{TR4}
\begin{enumerate}[label=\textbf{(TR\arabic*)}]
\item\label{TR1} \textbf{(Identité et complétion.)}
  \begin{enumerate}[label=(\alph*)]
  \item Pour tout objet $A$, le triangle
    $A \xrightarrow{\id} A \to 0 \to A[1]$ est distingué.
  \item Tout morphisme $u \colon A \to B$ peut être complété en un
    triangle distingué $A \xrightarrow{u} B \to C \to A[1]$.
  \item Tout triangle isomorphe à un triangle distingué est distingué.
  \end{enumerate}

\item\label{TR2} \textbf{(Rotation.)}
  Un triangle $A \xrightarrow{u} B \xrightarrow{v} C \xrightarrow{w} A[1]$
  est distingué si et seulement si
  $B \xrightarrow{v} C \xrightarrow{w} A[1] \xrightarrow{-u[1]} B[1]$
  est distingué.

\item\label{TR3} \textbf{(Morphisme de triangles.)}
  Étant donnés deux triangles distingués et des morphismes $f, g$
  formant un carré commutatif comme ci-dessous, il existe un morphisme $h$
  complétant le diagramme en un morphisme de triangles :
\[
\begin{tikzcd}
A \arrow[r, "u"] \arrow[d, "f"] & B \arrow[r, "v"] \arrow[d, "g"]
  & C \arrow[r, "w"] \arrow[d, dashed, "h"] & A{[1]} \arrow[d, "f{[1]}"] \\
A' \arrow[r, "u'"] & B' \arrow[r, "v'"] & C' \arrow[r, "w'"] & A'{[1]}
\end{tikzcd}
\]

\item\label{TR4} \textbf{(Axiome de l'octaèdre.)}
\index{axiome!de l'octaèdre}
  Étant donnés des morphismes $u \colon A \to B$ et $v \colon B \to C$,
  et des triangles distingués
  \begin{align*}
    & A \xrightarrow{u} B \to D \to A[1], \\
    & B \xrightarrow{v} C \to E \to B[1], \\
    & A \xrightarrow{v \circ u} C \to F \to A[1],
  \end{align*}
  il existe un triangle distingué $D \to F \to E \to D[1]$ tel que
  le diagramme complet commute.
\end{enumerate}
\end{definition}

\begin{theoreme}\label{thm:K-triangulated}
\index{K(A)@$K(\mathcal{A})$!structure triangulée}
\index{catégorie!homotopique!triangulée}
La catégorie homotopique $K(\mathcal{A})$, munie du foncteur de décalage
$[1]$ et des triangles de la forme
\[
  A^\bullet \xrightarrow{f} B^\bullet \xrightarrow{\iota}
  \operatorname{Cone}(f) \xrightarrow{\pi} A[1]^\bullet
\]
(et de leurs isomorphes), est une catégorie triangulée.
\end{theoreme}

\begin{proof}
Nous vérifions les axiomes.

\ref{TR1} : (a) Le cône de $\id \colon A^\bullet \to A^\bullet$ est
homotopiquement trivial : $\operatorname{Cone}(\id)^n = A^{n+1} \oplus A^n$
avec la différentielle qui en fait un complexe acyclique (exercice).
(b) Par définition, tout $f$ donne un triangle via le cône.
(c) Par définition de la classe $\Delta$.

\ref{TR2} : La rotation des triangles se vérifie par un calcul explicite
d'homotopie entre $\operatorname{Cone}(\iota)$ et $A[1]^\bullet$.

\ref{TR3} : On construit $h$ composante par composante en utilisant les
propriétés universelles dans la catégorie homotopique.

\ref{TR4} : Vérification par construction explicite, voir par exemple
Weibel \cite[1.5]{weibel1994}.
\end{proof}

\begin{corollaire}\label{cor:D-triangulated}
\index{D(A)@$D(\mathcal{A})$!structure triangulée}
La catégorie dérivée $D(\mathcal{A})$ hérite d'une structure triangulée.
De même pour $D^+(\mathcal{A})$, $D^-(\mathcal{A})$ et $D^b(\mathcal{A})$.
\end{corollaire}


%% ============================================================
\section{Foncteurs dérivés}\label{sec:derived-functors}
%% ============================================================

\begin{definition}[Foncteur dérivé à droite]\label{def:right-derived}
\index{foncteur!dérivé à droite|textbf}
\index{RF@$RF$}
Soit $F \colon \mathcal{A} \to \mathcal{B}$ un foncteur exact à gauche
entre catégories abéliennes, avec $\mathcal{A}$ ayant assez d'injectifs.
Le \textbf{foncteur dérivé à droite} $RF \colon D^+(\mathcal{A}) \to D^+(\mathcal{B})$
est défini comme suit :
\begin{enumerate}[label=\arabic*.]
\item Pour $A^\bullet \in D^+(\mathcal{A})$, on choisit une résolution
  injective $A^\bullet \xrightarrow{\sim} I^\bullet$ ;
\item On pose $RF(A^\bullet) = F(I^\bullet)$
  (application de $F$ degré par degré).
\end{enumerate}
Ceci est bien défini (indépendant du choix de résolution, à isomorphisme
canonique près dans $D^+(\mathcal{B})$) par le lemme de comparaison
(théorème~\ref{thm:comparison-lemma}).
\end{definition}

\begin{definition}[Foncteur dérivé à gauche]\label{def:left-derived}
\index{foncteur!dérivé à gauche|textbf}
\index{LF@$LF$}
Dualement, si $F \colon \mathcal{A} \to \mathcal{B}$ est exact à droite
et $\mathcal{A}$ a assez de projectifs, le \textbf{foncteur dérivé à gauche}
$LF \colon D^-(\mathcal{A}) \to D^-(\mathcal{B})$ est défini par
$LF(A^\bullet) = F(P^\bullet)$ où $P^\bullet \xrightarrow{\sim} A^\bullet$
est une résolution projective.
\end{definition}

\begin{definition}[Foncteurs dérivés classiques]\label{def:classical-derived}
\index{foncteur!dérivé!classique}
Pour un objet $A$ vu comme complexe concentré en degré $0$, on pose :
\[
  R^n F(A) = H^n(RF(A)) \quad \text{pour tout } n \geq 0.
\]
On a $R^0 F = F$ (car $F$ est exact à gauche). De même,
$L_n F(A) = H^{-n}(LF(A))$ avec $L_0 F = F$.
\end{definition}

\begin{definition}[$\operatorname{Ext}$ et $\operatorname{Tor}$]%
\label{def:ext-tor}
\index{Ext@$\operatorname{Ext}$|textbf}
\index{Tor@$\operatorname{Tor}$|textbf}
Soit $R$ un anneau et $A, B$ des $R$-modules.
\begin{enumerate}[label=(\roman*)]
\item Les \textbf{groupes d'extension}\index{groupe d'extension} sont
\[
  \operatorname{Ext}_R^n(A, B) = R^n \Hom_R(A, -)(B)
  = H^n(\Hom_R(A, I^\bullet))
\]
où $B \to I^\bullet$ est une résolution injective.
Alternativement,
$\operatorname{Ext}_R^n(A, B) = H^n(\Hom_R(P_\bullet, B))$
où $P_\bullet \to A$ est une résolution projective.

\item Les \textbf{groupes de torsion}\index{groupe de torsion!Tor} sont
\[
  \operatorname{Tor}_n^R(A, B) = L_n(A \otimes_R -)(B)
  = H_n(A \otimes_R P_\bullet)
\]
où $P_\bullet \to B$ est une résolution projective.
\end{enumerate}
\end{definition}

\begin{remarque}\label{rem:ext-tor-balanced}
\index{Ext@$\operatorname{Ext}$!balancé}
\index{Tor@$\operatorname{Tor}$!balancé}
Les foncteurs $\operatorname{Ext}$ et $\operatorname{Tor}$ sont
\emph{balancés} : on peut calculer $\operatorname{Ext}_R^n(A, B)$
soit via une résolution projective de $A$, soit via une résolution
injective de $B$, et on obtient le même résultat (à isomorphisme
canonique près). De même, $\operatorname{Tor}_n^R(A, B)$ peut être
calculé en résolvant soit $A$, soit $B$.
\end{remarque}

\begin{proposition}\label{prop:ext0-tor0}
\index{Ext@$\operatorname{Ext}$!en degré zéro}
\index{Tor@$\operatorname{Tor}$!en degré zéro}
On a les identifications :
\[
  \operatorname{Ext}_R^0(A, B) \cong \Hom_R(A, B), \qquad
  \operatorname{Tor}_0^R(A, B) \cong A \otimes_R B.
\]
\end{proposition}

\begin{proof}
Pour $\operatorname{Ext}^0$ : $H^0(\Hom_R(P_\bullet, B))
= \ker(\Hom(P_0, B) \to \Hom(P_1, B))$, qui s'identifie
à $\Hom(A, B)$ par exactitude de
$0 \to \Hom(A, B) \to \Hom(P_0, B) \to \Hom(P_1, B)$
(le foncteur $\Hom(-, B)$ est exact à gauche).
L'argument pour $\operatorname{Tor}_0$ est analogue.
\end{proof}


%% ============================================================
\section{Suite exacte longue en cohomologie}\label{sec:long-exact}
%% ============================================================

\begin{theoreme}[Suite exacte longue]\label{thm:long-exact-sequence}
\index{suite exacte!longue|textbf}
\index{morphisme connectant!suite exacte longue}
Soit $0 \to A^\bullet \xrightarrow{f} B^\bullet \xrightarrow{g} C^\bullet \to 0$
une suite exacte courte de complexes dans une catégorie abélienne. Alors
il existe des morphismes connectants
$\delta^n \colon H^n(C^\bullet) \to H^{n+1}(A^\bullet)$ et une suite
exacte longue
\[
\begin{tikzcd}[column sep=small]
\cdots \arrow[r]
& H^n(A^\bullet) \arrow[r, "H^n(f)"]
& H^n(B^\bullet) \arrow[r, "H^n(g)"]
& H^n(C^\bullet) \arrow[r, "\delta^n"]
& H^{n+1}(A^\bullet) \arrow[r]
& \cdots
\end{tikzcd}
\]
De plus, $\delta^n$ est naturel en les suites exactes courtes.
\end{theoreme}

\begin{proof}
\textbf{Construction de $\delta^n$.}
Soit $[c] \in H^n(C^\bullet)$, représenté par $c \in C^n$ avec
$d_C^n(c) = 0$. Comme $g^n$ est surjectif, il existe $b \in B^n$
avec $g^n(b) = c$. Considérons $d_B^n(b) \in B^{n+1}$.
On a
\[
  g^{n+1}(d_B^n(b)) = d_C^n(g^n(b)) = d_C^n(c) = 0,
\]
donc $d_B^n(b) \in \ker g^{n+1} = \im f^{n+1}$, et il existe un unique
$a \in A^{n+1}$ avec $f^{n+1}(a) = d_B^n(b)$.

\emph{$a$ est un cocycle.} On a $f^{n+2}(d_A^{n+1}(a)) = d_B^{n+1}(f^{n+1}(a))
= d_B^{n+1}(d_B^n(b)) = 0$, et comme $f^{n+2}$ est injectif,
$d_A^{n+1}(a) = 0$.

\emph{Bonne définition.} Si $b'$ est un autre relèvement de $c$, alors
$b - b' \in \ker g^n = \im f^n$, donc $b - b' = f^n(a_0)$ et
$d_B^n(b) - d_B^n(b') = d_B^n(f^n(a_0)) = f^{n+1}(d_A^n(a_0))$,
d'où $a - a' = d_A^n(a_0) \in \im d_A^n$.

On pose $\delta^n([c]) = [a] \in H^{n+1}(A^\bullet)$.

\medskip
\textbf{Exactitude en $H^n(C^\bullet)$.}
$\delta^n \circ H^n(g) = 0$ : si $[c] = H^n(g)([b_0])$ avec $d_B^n(b_0) = 0$,
on peut prendre $b = b_0$ et alors $d_B^n(b) = 0$, donc $a = 0$.

Réciproquement, si $\delta^n([c]) = 0$, alors $a = d_A^n(a_0)$ et
$b' = b - f^n(a_0)$ vérifie $d_B^n(b') = f^{n+1}(a - d_A^n(a_0)) = 0$ et
$g^n(b') = c$, donc $[c] = H^n(g)([b'])$.

\medskip
\textbf{Exactitude en $H^{n+1}(A^\bullet)$.}
$H^{n+1}(f) \circ \delta^n = 0$ : $H^{n+1}(f)([a]) = [f^{n+1}(a)]
= [d_B^n(b)] = 0$ dans $H^{n+1}(B^\bullet)$.

Réciproquement, si $[a] \in \ker H^{n+1}(f)$, alors $f^{n+1}(a) = d_B^n(b)$
pour un $b \in B^n$, et $c = g^n(b)$ est un cocycle (car
$d_C^n(c) = g^{n+1}(d_B^n(b)) = g^{n+1}(f^{n+1}(a)) = 0$), avec
$\delta^n([c]) = [a]$.

L'exactitude en $H^n(A^\bullet)$ et $H^n(B^\bullet)$ se vérifie de manière
analogue.
\end{proof}

\begin{corollaire}\label{cor:long-exact-derived}
\index{foncteur!dérivé!suite exacte longue}
\index{Ext@$\operatorname{Ext}$!suite exacte longue}
Soit $F \colon \mathcal{A} \to \mathcal{B}$ un foncteur exact à gauche
et $0 \to A \to B \to C \to 0$ une suite exacte courte dans $\mathcal{A}$.
Alors on a une suite exacte longue
\[
\begin{tikzcd}[column sep=small]
0 \arrow[r]
& FA \arrow[r]
& FB \arrow[r]
& FC \arrow[r, "\delta^0"]
& R^1 F(A) \arrow[r]
& R^1 F(B) \arrow[r]
& R^1 F(C) \arrow[r]
& \cdots
\end{tikzcd}
\]
En particulier, pour $F = \Hom_R(M, -)$ :
\[
\begin{tikzcd}[column sep=small]
0 \arrow[r]
& \Hom(M, A) \arrow[r]
& \Hom(M, B) \arrow[r]
& \Hom(M, C) \arrow[r]
& \operatorname{Ext}^1(M, A) \arrow[r]
& \cdots
\end{tikzcd}
\]
\end{corollaire}

\begin{proof}
On applique $RF$ à la suite exacte courte (vue comme triangle distingué
dans $D^+(\mathcal{A})$), ce qui donne un triangle distingué dans
$D^+(\mathcal{B})$. La suite exacte longue en cohomologie associée
à ce triangle donne le résultat.
\end{proof}


%% ============================================================
\section{Exercices}\label{sec:exos-derived}
%% ============================================================

\begin{exercice}\label{exo:cone-id-acyclic}
\index{cône!de mapping!de l'identité}
Montrer que $\operatorname{Cone}(\id_{A^\bullet})$ est un complexe acyclique
(c'est-à-dire $H^n(\operatorname{Cone}(\id)) = 0$ pour tout $n$),
et construire explicitement une homotopie contractante.
\end{exercice}

\begin{exercice}\label{exo:quasi-iso-cone}
\index{quasi-isomorphisme!et cône}
Montrer qu'un morphisme de complexes $f \colon A^\bullet \to B^\bullet$
est un quasi-isomorphisme si et seulement si
$\operatorname{Cone}(f)$ est acyclique.
\end{exercice}

\begin{exercice}\label{exo:ext-cyclic}
\index{Ext@$\operatorname{Ext}$!groupes cycliques}
Calculer $\operatorname{Ext}_\Z^n(\Z/m\Z, \Z/k\Z)$ pour tous $m, k \geq 1$
et tout $n \geq 0$.
\textit{Indication : utiliser la résolution projective
$0 \to \Z \xrightarrow{\times m} \Z \to \Z/m\Z \to 0$.}
\end{exercice}

\begin{exercice}\label{exo:tor-cyclic}
\index{Tor@$\operatorname{Tor}$!groupes cycliques}
Calculer $\operatorname{Tor}_n^\Z(\Z/m\Z, \Z/k\Z)$ pour tous $m, k \geq 1$
et tout $n \geq 0$.
\end{exercice}

\begin{exercice}\label{exo:ext-ses-interpretation}
\index{Ext@$\operatorname{Ext}$!interprétation par extensions}
\index{extension!de modules}
Soit $R$ un anneau et $A, C$ des $R$-modules. Montrer que
$\operatorname{Ext}_R^1(C, A)$ classifie les suites exactes courtes
$0 \to A \to B \to C \to 0$ à équivalence de Yoneda près.

\textit{Indication : à une telle suite, associer l'image de $\id_C$ par
le morphisme connectant $\Hom(C, C) \to \operatorname{Ext}^1(C, A)$.}
\end{exercice}

\begin{exercice}\label{exo:triangulated-abelian}
\index{catégorie!triangulée!pas abélienne}
Montrer qu'une catégorie triangulée non triviale n'est pas abélienne.

\textit{Indication : dans une catégorie triangulée, $\Hom(A, B)$ est
un groupe abélien et les suites exactes sont remplacées par des triangles.
Montrer que les biproduits ne coïncident pas avec les produits en général
dans une catégorie triangulée.}
\end{exercice}

\begin{exercice}\label{exo:derived-category-Z}
\index{D(Ab)@$D(\Ab)$}
\index{catégorie!dérivée!de $\Ab$}
Soit $A$ un groupe abélien. Montrer que dans $D(\Ab)$, $A$ (vu comme
complexe concentré en degré $0$) est isomorphe à toute résolution
projective ou injective de $A$.
\end{exercice}

\begin{exercice}\label{exo:horseshoe-derived}
\index{lemme!du fer à cheval!et foncteurs dérivés}
En utilisant le lemme du fer à cheval (exercice~\ref{exo:horseshoe}),
donner une preuve directe de l'existence de la suite exacte longue
des foncteurs dérivés (corollaire~\ref{cor:long-exact-derived})
sans passer par les catégories dérivées.
\end{exercice}

\begin{exercice}\label{exo:dimension-homologique}
\index{dimension homologique}
\index{dimension projective}
La \textbf{dimension projective}\index{dimension projective|textbf} d'un
$R$-module $M$, notée $\operatorname{pd}_R(M)$, est la longueur minimale
d'une résolution projective de $M$ (ou $+\infty$ si aucune n'est finie).
\begin{enumerate}[label=(\alph*)]
\item Montrer que $\operatorname{pd}_R(M) \leq n$ si et seulement si
  $\operatorname{Ext}_R^{n+1}(M, N) = 0$ pour tout $N$.
\item La \textbf{dimension globale}\index{dimension!globale} de $R$ est
  $\operatorname{gldim}(R) = \sup_M \operatorname{pd}_R(M)$.
  Montrer que $\operatorname{gldim}(\Z) = 1$.
\item Montrer que $\operatorname{gldim}(k) = 0$ pour tout corps $k$.
\end{enumerate}
\end{exercice}

\begin{exercice}\label{exo:rotation-triangle}
\index{triangle!distingué!rotation}
Vérifier en détail l'axiome \ref{TR2} pour la catégorie homotopique
$K(\mathcal{A})$ : si $A^\bullet \xrightarrow{f} B^\bullet \xrightarrow{\iota}
\operatorname{Cone}(f) \xrightarrow{\pi} A[1]^\bullet$ est un triangle
distingué, construire un isomorphisme explicite entre
$\operatorname{Cone}(\iota)$ et $A[1]^\bullet$ dans $K(\mathcal{A})$.
\end{exercice}

\begin{exercice}\label{exo:derived-adjunction}
\index{foncteur!dérivé!et adjonction}
Soit $(F, G)$ un couple de foncteurs adjoints entre catégories abéliennes
avec assez d'injectifs et de projectifs. Si $F$ est exact à droite et $G$
est exact à gauche, montrer que $(LF, RG)$ est un couple de foncteurs
adjoints au niveau des catégories dérivées :
\[
  \Hom_{D^-(\mathcal{B})}(LF(A^\bullet), B^\bullet)
  \cong \Hom_{D^-(\mathcal{A})}(A^\bullet, RG(B^\bullet)).
\]
\end{exercice}

\index{complexe de chaînes|)}
\index{catégorie!dérivée|)}
